\subsection{TLB 相关指令实现}

\begin{table}[h]
\centering
\caption{TLB指令执行特性对比}
\begin{tabular}{|c|c|c|c|}
\hline
\textbf{指令} & \textbf{执行级} & \textbf{冲突处理} & \textbf{备注} \\
\hline
TLBSRCH & EX & CSR读冲突用阻塞 & 查找结果在WB级写入CSR \\
\hline
TLBRD & WB & TLBIDX读冲突用阻塞 & 读出TLB项到CSR \\
\hline
TLBWR & WB & 写TLB用重取指 & 指定位置写入 \\
\hline
TLBFILL & WB & 写TLB用重取指 & 随机位置写入 \\
\hline
INVTLB & EX & 写TLB用重取指 & 有选择无效项 \\
\hline
\end{tabular}
\end{table}

\subsubsection{TLBSRCH指令的实现}

TLBSRCH指令用于查询TLB中是否存在匹配项。该指令的功能是：将CSR.ASID和CSR.TLBEHI的信息用于查询TLB，如果有命中项，将命中项的索引写入CSR.TLBIDX的Index域且NE位置0；否则将NE位置1。

TLBSRCH指令的执行涉及以下内容：

\begin{enumerate}
    \item \textbf{查找端口复用}：为避免新增大量查找逻辑，TLBSRCH指令复用访存指令的TLB查找端口，在EX级发起查找请求。
    
    % CSR读写冲突处理见流水线集成章节
    
    \item \textbf{访存虚地址来源切换}：TLBSRCH指令查找使用的VPPN来自CSR.TLBEHI，而load/store指令使用的VPPN来自计算得到的虚地址。在EX级需要根据指令类型选择合适的VPPN源。
\end{enumerate}

在stage\_ex.v中，使用ex\_inst\_tlbsrch信号识别TLBSRCH指令，并相应调整TLB查找输入：

\begin{lstlisting}[language=Verilog, caption=TLBSRCH虚地址源选择逻辑]
// For TLBSRCH, use CSR.TLBEHI; for load/store, use data_vaddr
assign s1_vppn = ex_inst_tlbsrch ? 
    u_csr.csr_tlbehi[`CSR_TLBEHI_VPPN] : 
    data_vaddr[31:13];
assign s1_va_bit12 = ex_inst_tlbsrch ? 1'b0 : data_vaddr[12];
\end{lstlisting}

% CSR依赖性检查见流水线集成章节

执行流程：

\begin{enumerate}
    \item \textbf{ID级}：识别TLBSRCH指令，检查TLBEHI和ASID依赖
    \item \textbf{EX级}：
    \begin{itemize}
        \item 从CSR.TLBEHI和CSR.ASID读取查找参数
        \item 将VPPN和ASID发送到TLB查找模块
        \item 获得查找结果（found和index）
    \end{itemize}
    \item \textbf{WB级}：
    \begin{itemize}
        \item 更新CSR.TLBIDX的Index域（查找命中时）
        \item 更新CSR.TLBIDX的NE位（命中为0，未命中为1）
    \end{itemize}
\end{enumerate}

\subsubsection{TLBRD指令的实现}

TLBRD指令用于读取TLB表项内容。该指令以CSR.TLBIDX的Index域值作为索引读取TLB表项，并将读出的数据填入相关CSR（TLBEHI、TLBELO0、TLBELO1、TLBIDX的PS域）。

\begin{enumerate}
    \item \textbf{读取时序}：TLBRD指令在WB级进行读取操作，确保读到的CSR.TLBIDX值和写入的结果都是正确的。
    
    \item \textbf{索引来源}：TLBRD指令读取CSR.TLBIDX的Index域作为TLB读地址。该CSR值可能被前面的CSRWR、CSRXCHG或TLBSRCH指令修改。
    
    % NE位处理见CSR寄存器章节
\end{enumerate}

在csr.v中的TLBRD执行逻辑：

\begin{lstlisting}[language=Verilog, caption=TLBRD执行逻辑]
// TLBRD: Read TLB entry based on Index field
else if (tlbrd_en) begin
    // Update TLBIDX: NE bit and PS field
    csr_tlbidx[`CSR_TLBIDX_PS] <= r_e ? r_ps : 6'b0;
    csr_tlbidx[`CSR_TLBIDX_NE] <= ~r_e;
    
    // Update TLBEHI with VPPN and ASID
    csr_tlbehi[`CSR_TLBEHI_VPPN] <= r_vppn;
    
    // Update TLBELO0 and TLBELO1
    csr_tlbelo0[`CSR_TLBELO_PPN]  <= r_ppn0;
    csr_tlbelo0[`CSR_TLBELO_PLV]  <= r_plv0;
    csr_tlbelo0[`CSR_TLBELO_MAT]  <= r_mat0;
    csr_tlbelo0[`CSR_TLBELO_D]    <= r_d0;
    csr_tlbelo0[`CSR_TLBELO_V]    <= r_v0;
    
    csr_tlbelo1[`CSR_TLBELO_PPN]  <= r_ppn1;
    csr_tlbelo1[`CSR_TLBELO_PLV]  <= r_plv1;
    csr_tlbelo1[`CSR_TLBELO_MAT]  <= r_mat1;
    csr_tlbelo1[`CSR_TLBELO_D]    <= r_d1;
    csr_tlbelo1[`CSR_TLBELO_V]    <= r_v1;
    
    // Update ASID
    csr_asid[`CSR_ASID_ASID] <= r_asid;
end
\end{lstlisting}

% 相关修复见CSR寄存器章节

\noindent\textbf{冲突处理}

TLBRD指令更新CSR会引发两类冲突：

\begin{enumerate}
    % 写后读冲突与重取指机制见流水线集成章节
\end{enumerate}

\subsubsection{TLBWR指令的实现}

TLBWR指令将CSR中的TLB表项信息写入TLB的指定项。被写入的表项信息来自CSR.TLBEHI、CSR.TLBELO0、CSR.TLBELO1和CSR.TLBIDX.PS，写入位置由CSR.TLBIDX的Index域指定。

\begin{enumerate}
    \item \textbf{写入时序}：TLBWR指令在WB级进行写入操作，确保读取到的所有CSR值都是正确的。
    
    \item \textbf{表项有效性判断}：若此时CSR.ESTAT.Ecode=0x3F（处于TLB重填异常处理中），则填入有效项；否则根据CSR.TLBIDX.NE位判断：NE=1时填入无效项，NE=0时填入有效项。
    
    % 写后读冲突与重取指机制见流水线集成章节
\end{enumerate}

在csr.v中的TLBWR执行逻辑：

\begin{lstlisting}[language=Verilog, caption=TLBWR表项有效性处理]
// TLBWR: Write TLB entry
wire tlbwr_en = /* TLBWR instruction detected */;
wire w_e_value;

// Determine entry validity based on exception state
if (wb_ecode == `ECODE_TLBR) begin
    // In TLB refill exception handler: always write valid entry
    w_e_value = 1'b1;
end else begin
    // Normal case: use NE bit to determine validity
    w_e_value = ~csr_tlbidx[`CSR_TLBIDX_NE];
end
\end{lstlisting}

% 重取指机制见流水线集成章节

\subsubsection{TLBFILL指令的实现}

TLBFILL指令功能与TLBWR类似，都是将TLB相关CSR中的页表信息写入TLB。区别在于：TLBWR指定写入位置（通过Index），而TLBFILL由硬件随机选择写入位置。

\begin{enumerate}
    \item \textbf{随机索引生成}：TLBFILL需要一个随机索引来决定写入TLB的哪一项。实现中使用一个简单的线性反馈移位寄存器(LFSR)或计数器生成伪随机数。
    
    \item \textbf{写入时序}：与TLBWR相同，在WB级进行写入操作。
    
    % 异常处理特殊情况见异常处理章节
\end{enumerate}

在stage\_ex.v中生成随机索引：

\begin{lstlisting}[language=Verilog, caption=TLBFILL随机索引生成]
// Random index generation for TLBFILL
reg [3:0] random_index;

// Simple LFSR-based random generator
always @(posedge clk) begin
    if (rst) begin
        random_index <= 4'b0001;
    end else begin
        random_index <= {random_index[2:0], 
                        random_index[3] ^ random_index[2]};
    end
end

// Pass random index to CSR module when TLBFILL executes
wire tlbfill_index = inst_tlbfill ? random_index : 4'b0;
\end{lstlisting}

\subsubsection{INVTLB指令的实现}

INVTLB指令用于无效TLB中的表项，维持TLB与内存页表的一致性。该指令支持多种操作类型（op=0~6），可以有选择地无效某些TLB表项。

\begin{enumerate}
    \item \textbf{并行匹配机制}：INVTLB指令需要对所有TLB表项进行匹配，判断哪些项需要无效。为支持高效执行，使用并行匹配逻辑，同时对所有表项进行匹配判断。
    
    \item \textbf{子匹配分解}：将各操作类型的匹配条件分解为多个"子匹配"：
    \begin{itemize}
        \item cond1：G位是否为0
        \item cond2：G位是否为1  
        \item cond3：源操作数指定的ASID是否匹配
        \item cond4：源操作数指定的VA是否匹配VPPN和PS
    \end{itemize}
    
    \item \textbf{操作类型编码}：
    \begin{itemize}
        \item op=0,1：无效所有项
        \item op=2：无效所有G=1的项
        \item op=3：保留（不执行）
        \item op=4：无效所有G=0且ASID匹配的项
        \item op=5：无效G=0、ASID匹配、VA匹配的项
        \item op=6：无效G=1或(ASID匹配且VA匹配)的项
    \end{itemize}
\end{enumerate}

在stage\_id.v中的INVTLB指令解码和有效性检查：

\begin{lstlisting}[language=Verilog, caption=INVTLB指令解码]
// Extract INVTLB operation code from instruction bits [4:0]
wire [4:0] invtlb_op = inst[4:0];

// Check if operation code is valid
wire invtlb_op_valid = (invtlb_op == 5'h0) |
                       (invtlb_op == 5'h1) |
                       (invtlb_op == 5'h2) |
                       (invtlb_op == 5'h3) |
                       (invtlb_op == 5'h4) |
                       (invtlb_op == 5'h5) |
                       (invtlb_op == 5'h6);

// Valid INVTLB instruction only if opcode is valid
assign inst_invtlb = inst_invtlb_raw & invtlb_op_valid;
\end{lstlisting}

在TLB模块中的INVTLB执行逻辑：

\begin{lstlisting}[language=Verilog, caption=TLB INVTLB匹配逻辑]
// For each TLB entry i:
wire cond1_i = ~tlb_g[i];              // G=0
wire cond2_i = tlb_g[i];               // G=1
wire cond3_i = (invtlb_asid == tlb_asid[i]); // ASID match
wire cond4_i = (s1_vppn[...] == tlb_vppn[i][...]) & 
               ps_match; // VA match

// Match conditions for each operation
wire inv_match[7:0];
assign inv_match[0] = 1'b1;  // op=0: all
assign inv_match[1] = 1'b1;  // op=1: all
assign inv_match[2] = cond2_i;  // op=2: G=1
assign inv_match[3] = 1'b0;  // op=3: reserved
assign inv_match[4] = cond1_i & cond3_i;  // op=4: G=0 & ASID
assign inv_match[5] = cond1_i & cond3_i & cond4_i;  // op=5: full
assign inv_match[6] = (cond2_i | (cond3_i & cond4_i)); // op=6

// Invalidate matching entries
always @(posedge clk) begin
    if (invtlb_valid) begin
        case (invtlb_op[2:0])
            // ... handle each operation type
            default: ; // Invalid op, no action
        endcase
        
        for (i = 0; i < TLBNUM; i = i + 1) begin
            if (inv_match[i]) begin
                tlb_e[i] <= 1'b0;
            end
        end
    end
end
\end{lstlisting}

在mycpu\_top.v中处理INVTLB的ASID和虚地址信息传递：

\begin{lstlisting}[language=Verilog, caption=INVTLB信息流传递]
// Pass ASID and VA information through pipeline
wire [ 9:0] wb_invtlb_asid;
wire [31:0] wb_invtlb_vaddr;

// Connect to TLB module
tlb_instance(
    .invtlb_op(wb_invtlb_op),
    .invtlb_asid(wb_invtlb_asid),
    .invtlb_vppn(wb_invtlb_vaddr[31:13]),
    // ... other connections
);
\end{lstlisting}

% INVTLB相关写后读冲突与重取指机制见流水线集成章节
